\begin{abstract}
This project explores critical slowing down in the two-dimensional Ising model near the critical temperature $T_c$. Using Markov chain Monte Carlo methods, we implement both the Metropolis and Wolff algorithms to generate spin configurations and measure autocorrelation times of magnetization and energy. Magnetic susceptibility is determined with an external field B, and, in the absence of the field, via the fluctuation-dissipation theorem (FDT). The scaling of relaxation times allows estimation of the dynamic critical exponent z for each algorithm. Comparisons between local and cluster updates highlight their relative efficiency at criticality, showing how algorithmic choices affect the sampling of Markov chains in systems exhibiting critical slowing down and field-induced responses.
\end{abstract}